\section{Evaluation}
\label{sec:evaluation}

\subsection{Metrics}

\begin{itemize}
\item \textbf{Rank IC} (Spearman correlation across stocks between forecasts and realised market-neutral 5-day
  returns, per test date), reported as its mean and ICIR; t-statistics with \textbf{Hansen--Hodrick} standard
  errors at lag 4. \dec{2026-09-26, audit E-2}
\item \textbf{Negative log-likelihood}, three ways (\texttt{NLL\_single}, \texttt{NLL\_base},
  \texttt{NLL\_full}), with the volatility-only gain next to the correction gain.
\item \textbf{A horizon-consistent long-short book}, equal-weighted and dollar-neutral; non-overlapping or
  staggered. \open{default scheme}
\item \textbf{Stability across seeds}: dispersion of rank IC and the count of seeds below a stated IC threshold.
\item \textbf{Regime-conditional results}: by regime and in transition windows.
\end{itemize}

\begin{whybox}
\begin{itemize}
\item \textbf{Rank IC} measures what the thesis is about, the ordering of stocks, and is robust to outliers
  and to the scale of the forecasts. It is the standard metric for cross-sectional signals and is invariant
  to the market-neutral demeaning.
\item \textbf{Hansen--Hodrick}: daily 5-day targets share four days, so consecutive ICs are correlated and naive
  t-statistics are too large. Hansen--Hodrick corrects exactly for overlap of known length, and in our
  simulation it had better size than Newey--West.
\item \textbf{Three NLLs}: a mixture can lower the NLL just by giving stressed regimes a wider distribution,
  without forecasting returns any better. Separating the variance gain from the correction gain stops that
  being mistaken for predictive skill.
\item \textbf{The long-short book} translates IC into an economic quantity. Equal-weighted and dollar-neutral, it
  matches the target's definition; horizon-consistent, so overlapping signals are not double-counted.
\item \textbf{Stability}: in the closest precedent, the largest differences between designs were in how many
  seeds failed, not in average accuracy. Their failure criterion is volatility-specific, so a returns
  version is defined in advance.
\item \textbf{Regime-conditional results}: the theory predicts gains where the regime matters (stress and
  transitions). The precedent's gain was concentrated in sustained high volatility and was weak at
  transitions; a regime process with persistence and transition dynamics is aimed at exactly that gap.
\end{itemize}
\end{whybox}

\subsection{Reporting additions (decided 2026-10-09)}

\dec{2026-10-09, Q9 update} Four additions to every arm's evaluation (gates, control, base, baselines), adopted after
reviewing the LightGBM pipeline in \texttt{OP model/lgbm}. They are reporting only and were fixed before any
test-period number was computed. \impl{brief 12}
\begin{enumerate}
\item \textbf{Long leg against short leg}: the long-short spread split into the long leg (mean market-neutral return of
  the top quantile) and the short leg (minus the mean of the bottom quantile), plus the rank IC computed separately
  within the upper and lower halves of the forecast distribution.
\item \textbf{Decile monotonicity}: mean return by forecast decile, and the monotonic relation test of Patton \&
  Timmermann (2010), with a stationary bootstrap.
\item \textbf{By calendar year and named episode}: 2010--2024 year by year, and eight episodes fixed in advance:
  euro debt and flash crash 2010; US downgrade 2011; China devaluation and oil 2015--16; the 2018 fourth-quarter
  sell-off; the COVID crash; the 2020 rebound; the 2022 bear market; the 2023 regional-bank stress.
\item \textbf{Execution lag}: a setting \texttt{execution\_lag}; with lag $L$ the forecast made at the close of $t$ is
  scored against the return over $(t+L,\,t+L+h]$. Evaluation only; every headline table at $L=0$ and $L=1$.
\end{enumerate}

\begin{whybox}
\begin{itemize}
\item \textbf{Legs}: in the LightGBM pipeline nearly all the edge sat in the short leg (IC 0.028 against 0.004). If the
  regime gain lives only in the short leg during stress, that changes how it could be used (shorting is costlier and
  constrained).
\item \textbf{Monotonicity}: a good spread can come from one extreme decile. A forecast that orders stocks well should
  raise returns steadily from decile 1 to 10.
\item \textbf{Years and episodes}: one headline number can hide an edge that lived in one year. Fixing the episode
  dates now means they cannot be chosen to flatter a gate.
\item \textbf{Execution lag}: the book assumes trading at the close at which the signal is computed. The IC decay curve
  showed several signals carry most of their information in the first two or three days, so a one-day delay is the
  honest check.
\end{itemize}
\end{whybox}

\paragraph{Turnover without trading.} Nothing is traded, but the evaluation builds a paper long-short book every
rebalance, so its turnover, $\tfrac12\sum_i|w_{i,t}-w_{i,t-1}|$, measures how stable the forecasts' rankings are (0: same
names; 1: whole book replaced). A gate that switches regime moves every stock's expert weights at once, so its
forecasts jump and turnover spikes; this makes turnover an assumption-free measure of how usable each gate's forecasts
are. A cost-adjusted return additionally needs an assumed cost per unit traded, which is why it stays with
\open{Q9 (c)}.

\open{Q9} Which metric family is primary: forecast accuracy or portfolio performance. In Gu, Kelly \& Xiu
the two families favour different depths, so the primary family must be fixed in advance.

\subsection{Diagnostics of what the gate contributes}

\open{Q10} The frozen base buys identifiability of the correction term, not of what it contains. Two
failure modes: the base underfits and the experts relearn general signal; or all experts converge to the
same correction and the gate is decorative. Planned diagnostics:
\begin{enumerate}
\item \textbf{Decomposition} of the correction into a gate-independent part (mean expert output) and a
  gate-conditional part, with each part's share of explained variation.
\item \textbf{Gate permutation test}: replace the fitted gate weights by uniform weights or weights from random
  other dates and build a null distribution.
\item \textbf{Expert-weight diagnostic} and the \textbf{effective sample per expert}.
\item A \textbf{variance decomposition} of gate weights across dates and across industries within a date (now
  meaningful with the industry-level gate).
\end{enumerate}

\begin{whybox}[Why these diagnostics]
A good number for the mixture is not evidence that \emph{regimes} helped: the experts might simply be fixing
what the base missed, identically in every regime. The decomposition and the permutation test answer the
question the thesis actually asks, ``does it matter \emph{which} regime the gate says we are in?''. The
permutation test needs no extra data and is the strongest single piece of evidence available.
\end{whybox}

\subsection{Baselines}

\open{Q12} Candidates: a capacity-matched single network; a network matched to one expert; a penalised
linear model; gradient-boosted trees (e.g.\ LightGBM); the no-gate mixture; jointly trained priors.

\begin{whybox}[Why these candidates]
Every claim is relative to a baseline. The capacity-matched network separates ``the gate matters'' from
``more parameters help''; the linear model shows the nonlinear machinery does anything at all; trees answer
the predictable referee question (trees are strong on tabular data, and Gu, Kelly \& Xiu cannot separate
networks from trees statistically); the no-gate mixture is the cleanest measure of gating itself.
\end{whybox}

\subsection{Multiple testing and pre-registration}

\open{Q15} The test set may be reused across models; what biases results is \emph{using test performance to
choose}. Seeds are not multiple testing; configurations are (about 18 gate-by-depth cells plus construction
arms). Open: the formal correction (e.g.\ Benjamini--Hochberg) and the pre-registration document.

\begin{whybox}[Why this matters so much here]
The differences between gates are likely small. Picking the best of 18 noisy estimates and reporting it as
if chosen in advance can create a ``winner'' out of noise alone. Locking every setting in the pilot,
nominating the primary cell in advance and correcting for the number of comparisons is what makes the
final ranking of gates believable.
\end{whybox}
